% Primary-operation input for maps-owned D.2.48; no separate theorem number.
\subsubsection*{The full Pontryagin cube of the weighted Engel torsion class}
We use the actual class $u\in H^2(X_d;\mathbb Z/d)$ whose
restrictions to the normalized cusp components are specified
in D.2.48. In particular its degree on every scalar-fixed
rational curve $C_j$, $j\in\mathbb Z/d$, is one, and its
component restriction is represented by the integral toric
support
\[
                  \psi(r,s)=-\frac{r(r-1)}2+3s(s-1).
                                                               \tag{D.TP1}
\]
The calculation below uses these actual restrictions. Merely
choosing an arbitrary degree-two class that is a unit on
torsion would not fix its pure coordinate in the presence
of the free summand.

The scalar circle action is semifree. On a finite cap of
order $m_i$ its gamma translation has primitive numerator
$n_i$, with $\gcd(m_i,n_i)=1$. A nontrivial root power cannot
be compensated by the scalar direction, because that would
require $j n_i/m_i$ integral. Thus these caps are scalar-free.
The cusp has exactly the $d$ fixed curves $C_j$; their normal
bundles are $\mathcal O(-1)\oplus\mathcal O(-1)$ with weights
$(1,-1)$, and the other cusp orbits are free. These are the
marked geometric inputs of the weighted family.

Because $X_d$ is simply connected, $u$ has an equivariant
lift $\widehat u\in H^2_{S^1}(X_d;\mathbb Z/d)$.
In the Borel spectral sequence the only possible
differentials out of ordinary degree two have targets
involving $H^1(X_d;\mathbb Z/d)$ or $H^3(BS^1;\mathbb Z/d)$,
and both vanish. Write
\[
         \widehat u|_{C_j}=x+\lambda_j t,
         \qquad x=[C_j]^*,\quad t\in H^2(BS^1;\mathbb Z/d).
                                                               \tag{D.TP2}
\]
The difference of any two equivariant lifts of a fixed
ordinary class on the simply connected normalized
$dP_6$ is a multiple of $t$. Hence the differences of
the characters on its two scalar-fixed curves can be
read from any toric equivariant lift of (D.TP1).
The fixed edges are the translates of the edge from
$(0,j)$ to $(0,j+1)$; the support slope in their scalar
direction is
\[
               \psi(0,j+1)-\psi(0,j)=6j.
                                                               \tag{D.TP3}
\]
The fibre character of the associated monomial
trivialization is the negative of this slope, up to
one global character convention. Equivalently the
intervening weight-one orbit curve has degree $-6$.
Thus, for one common sign $\epsilon=\pm1$,
\[
                    \lambda_j=\lambda_0+6\epsilon j
                                 \quad\text{in }\mathbb Z/d.
                                                               \tag{D.TP4}
\]
This equation is compatible with the period $j\mapsto j+d$.
It does not assert that the integral support line itself
descends with an integral global circle linearization:
the needed character equality is modulo $d$ and follows
from the existing global equivariant cohomology lift.

Suppose now that $3\mid d$. Thomas's operation is defined
with precisely the required full coefficients:
\[
 \mathcal P_3:H^2(-;\mathbb Z/d)
                    \longrightarrow H^6(-;\mathbb Z/(3d)).
                                                               \tag{D.TP5}
\]
This is \cite[equation (2.1), p.267]{ffThomas56} with
$p=3$ and $m=d/3$. On the reduction of an integral
class $z$, it equals the reduction of $z^3$ modulo
$3d$. Indeed the correction term in Thomas's actual
cochain formula (3.4) is zero for an integral cocycle,
and its first term is the ordinary threefold diagonal
product \cite[equations (3.4)--(3.6), p.268]{ffThomas56}.
Each restriction (D.TP2) has such an integral lift.

Apply equivariant localization with coefficients
$\mathbb Z/(3d)$ to $\mathcal P_3(\widehat u)$.
No field assumption is needed here: the complement
is scalar-free, and each normal Euler class becomes
a unit after inverting $t$. On a fixed curve it is
\[
                    (-x+t)(-x-t)=-t^2,
\]
because $x^2=0$. Integrating the term of ordinary
curve degree two gives
\[
 \left\langle\mathcal P_3(u),[X_d]\right\rangle
                      =-3\sum_{j=0}^{d-1}\lambda_j^2
                                  \quad\text{in }\mathbb Z/(3d).
                                                               \tag{D.TP6}
\]
The right side is independent of the integer
representatives of the characters: $3\mid d$ makes
the cube of an integral lift well defined modulo
$3d$. The base cohomology polynomial ring injects
into its localization at $t$, so the degree-zero
identity is an identity in the original coefficient
ring, not just a localized statement.

For arbitrary integers $d>0,\lambda_0$ and either
sign in (D.TP4), the exact identity is
\[
 \sum_{j=0}^{d-1}(\lambda_0+6\epsilon j)^2
   =d\left(\lambda_0^2+6\epsilon\lambda_0(d-1)
                         +6(d-1)(2d-1)\right).              \tag{D.TP7}
\]
Substitution into (D.TP6) proves
\[
                 \left\langle\mathcal P_3(u),[X_d]\right\rangle=0
                                  \quad\text{in }\mathbb Z/(3d).
                                                               \tag{D.TP8}
\]
Thus the complete pure cyclic fundamental-class coordinate
is zero, including the order-three subgroup invisible to
the ordinary mod-$d$ cube. For $3\nmid d$ the same
localization calculation with the ordinary cube gives
zero modulo $d$. Mixed coordinates are separate from
this calculation; in particular it does not discard
the order-$2d$ mixed Pontryagin-square terms when $d$
is even.

The universal integer identities, support slope,
single-curve Euler/cube calculation and an explicitly
bounded supplementary check are executed in
\path{code/fourfolds_D48_thomas.py}.
The supplementary range is $d=3,6,\ldots,120$,
$\lambda_0=0,\ldots,d-1$ and both signs, giving
$4920$ exact checks. The arbitrary-$d$ conclusion
comes from (D.TP7), not that finite range.
